% from ega4-2-fr.tex
\oldpage[IV-2]{50}

un $\mathcal O_X$-Module cohérent, $y$ un point de $Y$~; supposons que $t_y$
soit $\mathcal F_y$-régulier et que $\mathcal F/t\mathcal F$ soit intègre au
point $y$ (3.2.4). Alors $\mathcal F$ est intègre au point $y$.
\end{proposition}

\begin{proof}
Compte tenu de (3.4.4), il suffit de prouver que $y$ est contenu dans une
\emph{seule} composante irréductible de $\operatorname{Supp}(\mathcal F)$, et
que si $s$ est le point générique de cette composante, on a
$\operatorname{long}(\mathcal F_s)=1$. Or, par hypothèse, $y$ n'appartient
qu'à une seule composante irréductible de
$\operatorname{Supp}(\mathcal F/t\mathcal F)$, et si $z$ est le point
générique de cette composante, on a
$\operatorname{long}((\mathcal F/t\mathcal F)_z)=1$~; la conclusion résulte
donc de (3.4.2).
\end{proof}

\begin{proposition}[3.4.6]
\label{IV.2.3.4.6-fr}
Les hypothèses étant celles de (3.4.1), soit $x$ un point de $Y$. Supposons
que $Y$ ne contienne aucun des $S_i$ contenant $x$, et que
$\mathcal F/t\mathcal F$ soit réduit au point $x$ (3.2.2). Alors $t_x$ est
$\mathcal F_x$-régulier et $\mathcal F$ est réduit au point $x$. En outre, si
$z_j$ est point générique d'une composante irréductible de
$\operatorname{Supp}(\mathcal F/t\mathcal F)$ contenant $x$, $z_j$ est
contenu dans un seul des $S_i$, et $\mathcal O_{S_i,z_j}$ est un anneau de
valuation discrète dont $t_{z_j}$ est une uniformisante.
\end{proposition}

\begin{proof}
Le fait que $t_x$ est $\mathcal F_x$-régulier résulte du lemme suivant appliqué
à l'anneau $\mathcal O_x$~:

\begin{lemma}[3.4.6.1]
\label{IV.2.3.4.6.1-fr}
Soient $A$ un anneau noethérien, $M$ un $A$-module de type fini,
$\mathfrak p_i$ les éléments minimaux de $\operatorname{Supp}(M)$, $t$ un
élément de $A$. On suppose que $t$ n'appartient à aucun des $\mathfrak p_i$,
et que $M/tM$ est un $A$-module réduit (3.2.2). Alors $t$ est $M$-régulier.
\end{lemma}

\begin{proof}
Tout idéal premier $\mathfrak p\in\operatorname{Supp}(M)$ contient un des
$\mathfrak p_i$~; comme $t$ n'appartient à aucun des $\mathfrak p_i$,
l'homothétie de rapport $t$ dans $M_{\mathfrak p}$ n'est pas nilpotente
(Bourbaki, \emph{Alg. comm.}, chap.~IV, \S~1, n\textsuperscript{o}~4, cor. de
la prop.~9). Désignons par $N$ le sous-module de $M$ formé des éléments
annulés par une puissance de $t$, et posons $P=M/N$~; nous allons montrer que
$N=0$. Comme $t$ est $P$-régulier, on a une suite exacte (3.4.1.4)
\[
0\longrightarrow N/tN\longrightarrow M/tM\longrightarrow P/tP
\longrightarrow0.
\]
Comme $N$ est de type fini, il est annulé par une puissance de $t$, et il
suffit donc de montrer que $N/tN=0$. Comme $N/tN$ est un sous-module de
$M/tM$, il suffit de prouver que $(N/tN)_{\mathfrak p}=0$ pour tout
$\mathfrak p\in\operatorname{Ass}(M/tM)$, ou encore que l'homomorphisme
\[
u_{\mathfrak p}:(M/tM)_{\mathfrak p}\longrightarrow(P/tP)_{\mathfrak p}
\]
est bijectif pour tout $\mathfrak p\in\operatorname{Ass}(M/tM)$. Or on a
$(P/tP)_{\mathfrak p}\neq0$~; en effet, comme
$\mathfrak p\in\operatorname{Supp}(M/tM)
=\operatorname{Supp}(M)\cap V(t)$ $(0_{\mathrm I},1.7.5)$, l'image de $t$ dans
$A_{\mathfrak p}$ est contenue dans l'idéal maximal
$\mathfrak pA_{\mathfrak p}$, donc l'hypothèse
$P_{\mathfrak p}=tP_{\mathfrak p}$ entraînerait $P_{\mathfrak p}=0$ par le
lemme de Nakayama~; on aurait donc $M_{\mathfrak p}=N_{\mathfrak p}$ et
l'homothétie de rapport $t$ dans $M_{\mathfrak p}$ serait nilpotente~; mais
cela contredit la remarque faite au début, puisque
$\mathfrak p\in\operatorname{Supp}(M)$. Cela étant, l'hypothèse que $M/tM$ est
réduit entraîne $\operatorname{long}((M/tM)_{\mathfrak p})=1$, et comme
$(P/tP)_{\mathfrak p}\neq0$, $u_{\mathfrak p}$ est nécessairement bijectif,
ce qui prouve le lemme.
\end{proof}

Par hypothèse, aucun des cycles premiers associés immergés de
$\mathcal F/t\mathcal F$ ne contient $x$, donc aucun des cycles premiers
associés immergés de $\mathcal F$ ne contient $x$, en vertu de (3.4.4).
D'autre part, appliquant (3.4.2) à une composante irréductible de
$\operatorname{Supp}(\mathcal F/t\mathcal F)$ contenant $x$, on voit que
$\operatorname{long}(\mathcal F_{s_i})=1$ pour tout $S_i$ contenant $x$, ce
qui achève de montrer que $\mathcal F_x$ est réduit~; enfin, les dernières
assertions sont aussi conséquences de (3.4.2).
\end{proof}

\begin{corollary}[3.4.7]
\label{IV.2.3.4.7-fr}
Soient $A$ un anneau local noethérien, $\mathfrak m$ son idéal maximal, $M$ un
$A$-module de type fini, $(x_i)_{1\leqslant i\leqslant k}$ une famille
d'éléments de $\mathfrak m$ formant une partie d'un système de

